\documentclass[11pt, a4paper]{article}

\usepackage[utf8]{inputenc}
\usepackage[T1]{fontenc}
\usepackage[ngerman]{babel}
\usepackage{amsmath}
\usepackage{amssymb}
\usepackage{graphicx}
\usepackage{geometry}
\usepackage{fancyhdr}
\usepackage{enumitem}
\usepackage{array}
\usepackage{eso-pic}

\usepackage{siunitx}
\usepackage[b]{esvect}
\DeclareSIUnit{\litre}{l}
\DeclareSIUnit{\barPr}{bar}

\sisetup{angle-symbol-degree = ^\circ}
\sisetup{locale = DE, separate-uncertainty, inter-unit-product = {},
range-units = brackets, list-units = single, per-mode=symbol-or-fraction}

\makeatletter
\def\input@path{{./}{../../}}
\makeatother
\graphicspath{{./}{../../}}

\geometry{a4paper, top=2.5cm, bottom=2.5cm, left=2.5cm, right=2.5cm}

\input{defs}

\pagestyle{fancy}
\fancyhf{}
\cfoot{\thepage}

\begin{document}
\noindent
\vspace{9cm}
\begin{center}
    {\Huge \textbf{Lösung}} \\[1.5em]
    {\Large \textbf{Physikalische Grundlagen}} \\[1em]
    {\large 4. Schriftliche Prüfung: Sommersemester 2026} \\[1em]
    {\large Studiengang: Clinical Engineering}
\end{center}

\newpage

\section*{Lösung Beispiel 1 -- Wasserkocher vergessen}
\begin{enumerate}
    \item \textbf{Temperatur-Zeit-Diagramm:}
    Die Temperatur steigt zunächst von \SI{20}{\degreeCelsius} auf \SI{100}{\degreeCelsius}. Danach bleibt sie bei \SI{100}{\degreeCelsius} konstant, während das Wasser verdampft.

    \item \textbf{Energie zum Erhitzen auf \SI{100}{\degreeCelsius}:}
    Zuerst berechnen wir die Masse des Wassers:
    \[
        m_W = \rho_W V = \SI{1}{\kilogram\per\deci\metre\cubed} \cdot \SI{1}{\litre} = \SI{1}{\kilogram}
    \]
    Für das Erwärmen gilt
    \[
        Q_1 = m_W c_W \Delta T
    \]
    mit
    \[
        \Delta T = \SI{100}{\degreeCelsius} - \SI{20}{\degreeCelsius} = \SI{80}{\kelvin}
    \]
    also
    \[
        Q_1 = \SI{1}{\kilogram} \cdot \SI{4.18}{\kilo\joule\per(\kilogram\cdot\kelvin)} \cdot \SI{80}{\kelvin} = \SI{334.4}{\kilo\joule}
    \]

    \item \textbf{Energie zum Verdampfen:}
    Es gilt
    \[
        Q_2 = m_W \lambda_{\text{verd}} = \SI{1}{\kilogram} \cdot \SI{2256}{\kilo\joule\per\kilogram} = \SI{2256}{\kilo\joule}
    \]

    \item \textbf{Energie des Wasserkochers in \SI{25}{\minute}:}
    Zuerst wird die Zeit in Sekunden umgerechnet:
    \[
        t = \SI{25}{\minute} = 25 \cdot 60\,\si{\second} = \SI{1500}{\second}
    \]
    Damit ergibt sich
    \[
        E_{\text{Kocher}} = P t = \SI{2000}{\watt} \cdot \SI{1500}{\second} = \SI{3000000}{\joule} = \SI{3000}{\kilo\joule}
    \]

    \item \textbf{Fazit:}
    Insgesamt benötigt man
    \[
        Q_{\text{ges}} = Q_1 + Q_2 = \SI{334.4}{\kilo\joule} + \SI{2256}{\kilo\joule} = \SI{2590.4}{\kilo\joule}
    \]
    Da
    \[
        E_{\text{Kocher}} = \SI{3000}{\kilo\joule} > Q_{\text{ges}} = \SI{2590.4}{\kilo\joule}
    \]
    kommen Sie \textbf{nicht} rechtzeitig zurück. Das gesamte Wasser ist bereits verdampft.
\end{enumerate}

\newpage

\section*{Lösung Beispiel 2 -- Schlittenfahren am Iglu}
\begin{enumerate}
    \item \textbf{Kräfte und Komponenten:}
    Es wirken die Gewichtskraft $\ivecS{F}{G}$ nach unten und die Normalkraft $\ivecS{F}{N}$ senkrecht zur Oberfläche nach außen.
    Für die Komponenten gilt:
    \[
        F_{G,r} = -mg\cos(\alpha), \qquad F_{G,t} = mg\sin(\alpha)
    \]
    \[
        F_{N,r} = F_N, \qquad F_{N,t} = 0
    \]

    \item \textbf{Bewegungsgleichungen:}
    In radialer Richtung:
    \[
        F_N - mg\cos(\alpha) = m a_r
    \]
    In tangentialer Richtung:
    \[
        mg\sin(\alpha) = m a_t
    \]

    \item \textbf{Bedingung für das Abheben:}
    Beim Abheben ist $F_N = 0$ und die radiale Beschleunigung ist die Zentripetalbeschleunigung
    \[
        a_r = -\frac{v^2}{R}
    \]
    Einsetzen in die radiale Gleichung ergibt
    \[
        0 - mg\cos(\alpha_{\text{max}}) = -m\frac{v^2}{R}
    \]
    also
    \[
        g\cos(\alpha_{\text{max}}) = \frac{v^2}{R}
    \]

    \item \textbf{Geschwindigkeit aus Energieerhaltung:}
    Die Abnahme der potenziellen Energie wird vollständig in kinetische Energie umgewandelt:
    \[
        mgR - mgR\cos(\alpha) = \frac{1}{2}mv^2
    \]
    Daraus folgt
    \[
        v^2 = 2gR(1-\cos(\alpha))
    \]

    \item \textbf{Bestimmung von $\alpha_{\text{max}}$:}
    Einsetzen von $v^2$ in die Abhebegleichung liefert
    \[
        g\cos(\alpha_{\text{max}}) = \frac{2gR(1-\cos(\alpha_{\text{max}}))}{R}
    \]
    also
    \[
        \cos(\alpha_{\text{max}}) = 2(1-\cos(\alpha_{\text{max}}))
    \]
    und damit
    \[
        3\cos(\alpha_{\text{max}}) = 2
    \]
    \[
        \cos(\alpha_{\text{max}}) = \frac{2}{3}
    \]
    Somit ergibt sich
    \[
        \alpha_{\text{max}} = \arccos\!\left(\frac{2}{3}\right) \approx \ang{48.2}
    \]
    Das Ergebnis ist unabhängig von der Masse $m$ und vom Radius $R$.
\end{enumerate}

\newpage

\section*{Lösung Beispiel 3 -- Reibungsfreies Schieben}
\begin{enumerate}
    \item \textbf{Freikörperbild für $m$:}
    Auf die kleine Masse wirken die Gewichtskraft $\ivecS{F}{G} = m\ivec{g}$ senkrecht nach unten und die Normalkraft $\ivecS{F}{N}$ senkrecht zur schiefen Ebene. Das ganze System bewegt sich mit der horizontalen Beschleunigung $a_x$.

    \item \textbf{Kräftegleichgewichte in $x$- und $y$-Richtung:}
    Wir zerlegen die Normalkraft in horizontale und vertikale Komponenten. Mit dem Neigungswinkel $\alpha$ gilt in vertikaler Richtung
    \[
        \sum F_y = F_N\cos(\alpha) - mg = 0
    \]
    also
    \[
        F_N\cos(\alpha) = mg
    \]
    In horizontaler Richtung gilt
    \[
        \sum F_x = F_N\sin(\alpha) = m a_x
    \]

    \item \textbf{Beschleunigung $a_x$:}
    Aus der vertikalen Gleichung folgt
    \[
        F_N = \frac{mg}{\cos(\alpha)}
    \]
    Einsetzen in die horizontale Gleichung ergibt
    \[
        \frac{mg}{\cos(\alpha)}\sin(\alpha) = m a_x
    \]
    also
    \[
        g\tan(\alpha) = a_x
    \]
    und damit
    \[
        a_x = g\tan(\alpha)
    \]

    \item \textbf{Gesamtkraft $F$:}
    Um das gesamte System der Masse $M+m$ mit dieser Beschleunigung zu bewegen, braucht man
    \[
        F = (M+m) a_x
    \]
    also
    \[
        F = (M+m) g\tan(\alpha)
    \]
\end{enumerate}

\newpage

\section*{Lösung Beispiel 4 -- Single-Choice}
\begin{enumerate}
    \item \textbf{Leistung}
    \item \textbf{Das Volumen}
    \item \textbf{Faktor 4}
    \item \textbf{dem Gewicht des verdrängten Mediums}
    \item \textbf{Kräfte treten paarweise als gleich große, entgegengesetzt gerichtete Wechselwirkungen auf.}
    \item \textbf{mittlere kinetische Energie der Teilchen}
    \item \textbf{Arbeit, Einheit: \si{\joule}}
    \item \textbf{Sein Wirkungsgrad hängt nur von den Temperaturen des heißen und kalten Reservoirs in Kelvin ab.}
    \item \textbf{Der Druck verdoppelt sich.}
    \item \textbf{Der zweite Hauptsatz}
\end{enumerate}

\end{document}
\documentclass[11pt, a4paper]{article}

\usepackage[utf8]{inputenc}
\usepackage[T1]{fontenc}
\usepackage[ngerman]{babel}
\usepackage{amsmath}
\usepackage{amssymb}
\usepackage{graphicx}
\usepackage{geometry}
\usepackage{fancyhdr}
\usepackage{enumitem}
\usepackage{array}
\usepackage{eso-pic}

\usepackage{siunitx}
\usepackage[b]{esvect}
\DeclareSIUnit{\litre}{l}
\DeclareSIUnit{\barPr}{bar}

\sisetup{angle-symbol-degree = ^\circ}
\sisetup{locale = DE, separate-uncertainty, inter-unit-product = {},
range-units = brackets, list-units = single, per-mode=symbol-or-fraction}

\makeatletter
\def\input@path{{./}{../../}}
\makeatother
\graphicspath{{./}{../../}}

\geometry{a4paper, top=2.5cm, bottom=2.5cm, left=2.5cm, right=2.5cm}

\input{defs}

\pagestyle{fancy}
\fancyhf{}
\cfoot{\thepage}

\begin{document}
\noindent
\vspace{9cm}
\begin{center}
    {\Huge \textbf{Lösung}} \\[1.5em]
    {\Large \textbf{Physikalische Grundlagen}} \\[1em]
    {\large 4. Schriftliche Prüfung: Sommersemester 2026} \\[1em]
    {\large Studiengang: Clinical Engineering}
\end{center}

\newpage

\section*{Lösung Beispiel 1 -- Wasserkocher vergessen}
\begin{enumerate}
    \item \textbf{Temperatur-Zeit-Diagramm:}
    Die Temperatur steigt zunächst von \SI{20}{\degreeCelsius} auf \SI{100}{\degreeCelsius}. Während des anschließenden Siedens bleibt die Temperatur bei \SI{100}{\degreeCelsius} konstant, obwohl weiterhin Energie zugeführt wird.

    \item \textbf{Energie zum Erhitzen auf \SI{100}{\degreeCelsius}:}
    Zunächst folgt aus Volumen und Dichte die Masse des Wassers:
    \[
        m_W = \rho_W V = \SI{1}{\kilogram\per\deci\metre\cubed} \cdot \SI{1}{\litre} = \SI{1}{\kilogram}
    \]
    Mit $\Delta T = \SI{100}{\degreeCelsius} - \SI{20}{\degreeCelsius} = \SI{80}{\kelvin}$ ergibt sich
    \[
        Q_1 = m_W c_W \Delta T = \SI{1}{\kilogram} \cdot \SI{4.18}{\kilo\joule\per(\kilogram\cdot\kelvin)} \cdot \SI{80}{\kelvin} = \SI{334.4}{\kilo\joule}
    \]

    \item \textbf{Energie zum Verdampfen:}
    \[
        Q_2 = m_W \lambda_{\text{verd}} = \SI{1}{\kilogram} \cdot \SI{2256}{\kilo\joule\per\kilogram} = \SI{2256}{\kilo\joule}
    \]

    \item \textbf{Vom Wasserkocher gelieferte Energie in \SI{25}{\minute}:}
    Die Zeit in SI-Einheiten lautet
    \[
        t = \SI{25}{\minute} = 25 \cdot 60\,\si{\second} = \SI{1500}{\second}
    \]
    Damit ist
    \[
        E_{\text{Kocher}} = Pt = \SI{2000}{\watt} \cdot \SI{1500}{\second} = \SI{3.0e6}{\joule} = \SI{3000}{\kilo\joule}
    \]

    \item \textbf{Fazit:}
    Insgesamt benötigt man
    \[
        Q_{\text{ges}} = Q_1 + Q_2 = \SI{334.4}{\kilo\joule} + \SI{2256}{\kilo\joule} = \SI{2590.4}{\kilo\joule}
    \]
    Da
    \[
        E_{\text{Kocher}} = \SI{3000}{\kilo\joule} > Q_{\text{ges}} = \SI{2590.4}{\kilo\joule}
    \]
    kommen Sie \textbf{nicht} rechtzeitig zurück. Das gesamte Wasser ist bereits verdampft.
\end{enumerate}

\newpage

\section*{Lösung Beispiel 2 -- Schlittenfahren am Iglu}
\begin{enumerate}
    \item \textbf{Kräfte und Komponenten:}
    Es wirken die Gewichtskraft $\ivecS{F}{G}$ nach unten und die Normalkraft $\ivecS{F}{N}$ senkrecht zur Oberfläche nach außen. Die Komponenten lauten
    \[
        F_{G,r} = -mg\cos(\alpha), \qquad F_{G,t} = mg\sin(\alpha)
    \]
    \[
        F_{N,r} = F_N, \qquad F_{N,t} = 0
    \]

    \item \textbf{Bewegungsgleichungen:}
    Für die radiale und tangentiale Richtung gilt
    \[
        F_N - mg\cos(\alpha) = m a_r
    \]
    \[
        mg\sin(\alpha) = m a_t
    \]

    \item \textbf{Bedingung für das Abheben:}
    Im Moment des Abhebens ist $F_N = 0$ und $a_r = -v^2/R$. Damit folgt
    \[
        0 - mg\cos(\alpha_{\text{max}}) = -m\frac{v^2}{R}
    \]
    also
    \[
        g\cos(\alpha_{\text{max}}) = \frac{v^2}{R}
    \]

    \item \textbf{Geschwindigkeit aus Energieerhaltung:}
    Die Abnahme der potenziellen Energie wird vollständig in kinetische Energie umgewandelt:
    \[
        mgR - mgR\cos(\alpha) = \frac{1}{2}mv^2
    \]
    Daraus erhält man
    \[
        v^2 = 2gR(1-\cos(\alpha))
    \]

    \item \textbf{Bestimmung von $\alpha_{\text{max}}$:}
    Einsetzen von $v^2$ in die Gleichung aus 3 liefert
    \[
        g\cos(\alpha_{\text{max}}) = \frac{2gR(1-\cos(\alpha_{\text{max}}))}{R}
    \]
    also
    \[
        \cos(\alpha_{\text{max}}) = 2(1-\cos(\alpha_{\text{max}}))
    \]
    \[
        3\cos(\alpha_{\text{max}}) = 2
    \]
    \[
        \cos(\alpha_{\text{max}}) = \frac{2}{3}
    \]
    und damit
    \[
        \alpha_{\text{max}} = \arccos\left(\frac{2}{3}\right) \approx \ang{48.2}
    \]
    Das Ergebnis ist unabhängig von der Masse $m$ und vom Radius $R$.
\end{enumerate}

\newpage

\section*{Lösung Beispiel 3 -- Reibungsfreies Schieben}
\begin{enumerate}
    \item \textbf{Freikörperbild für $m$:}
    Auf die kleine Masse wirken die Gewichtskraft $\ivecS{F}{G} = m\ivec{g}$ nach unten und die Normalkraft $\ivecS{F}{N}$ senkrecht zur schiefen Ebene. Das System wird horizontal mit der Beschleunigung $a_x$ bewegt.

    \item \textbf{Kräftegleichgewichte in $x$- und $y$-Richtung:}
    Wir zerlegen die Normalkraft. In vertikaler Richtung gilt wegen fehlender Vertikalbeschleunigung
    \[
        \sum F_y = F_N \cos(\alpha) - mg = 0
    \]
    also
    \[
        F_N \cos(\alpha) = mg
    \]
    In horizontaler Richtung gilt
    \[
        \sum F_x = F_N \sin(\alpha) = m a_x
    \]

    \item \textbf{Herleitung der Beschleunigung $a_x$:}
    Aus der vertikalen Gleichung folgt
    \[
        F_N = \frac{mg}{\cos(\alpha)}
    \]
    Eingesetzt in die horizontale Gleichung ergibt das
    \[
        \frac{mg}{\cos(\alpha)}\sin(\alpha) = m a_x
    \]
    und damit
    \[
        a_x = g\tan(\alpha)
    \]

    \item \textbf{Gesamtkraft auf das System:}
    Die Gesamtmasse $M+m$ muss mit dieser Beschleunigung bewegt werden. Daher
    \[
        F = (M+m)a_x = (M+m)g\tan(\alpha)
    \]
\end{enumerate}

\newpage

\section*{Lösung Beispiel 4 -- Single-Choice}
\begin{enumerate}
    \item \textbf{Leistung}
    \item \textbf{Das Volumen}
    \item \textbf{Faktor 4}
    \item \textbf{dem Gewicht des verdrängten Mediums}
    \item \textbf{Kräfte treten paarweise als gleich große, entgegengesetzt gerichtete Wechselwirkungen auf.}
    \item \textbf{mittlere kinetische Energie der Teilchen}
    \item \textbf{Arbeit, Einheit: \si{\joule}}
    \item \textbf{Sein Wirkungsgrad hängt nur von den Temperaturen des heißen und kalten Reservoirs in Kelvin ab.}
    \item \textbf{Der Druck verdoppelt sich.}
    \item \textbf{Der zweite Hauptsatz}
\end{enumerate}

\end{document}